% lemniro: {"kind":"note","description":"A complete proof that compact subsets of Hausdorff spaces are closed, with every use of the hypotheses made explicit.","topic":"Topology","series":"A first course in topology","order":2,"date":"2026-10-01","readingMinutes":14,"prerequisites":["Topological spaces","Open covers","The Hausdorff condition"],"related":["topological-spaces","finite-intersections"]}
\documentclass[11pt]{article}
\input{tex/lemniro-preamble}
\title{Why compact sets are closed in Hausdorff spaces}
\begin{document}
\maketitle

\section{From two points to a whole set}
The Hausdorff condition separates two distinct points by disjoint open
neighborhoods. Our goal is to separate a point from an entire compact
subset. Compactness makes that change of scale possible: it reduces an
open cover, which may be infinite, to finitely many of its members.

\begin{definition}[Compact subset]\label{def:compact}
A subset $K$ of a topological space $X$ is compact if every family
$\set{V_\alpha}_{\alpha\in A}$ of open subsets of $X$ with
$K\subseteq\bigcup_{\alpha\in A}V_\alpha$ has a finite subfamily
$V_{\alpha_1},\ldots,V_{\alpha_n}$ whose union still contains $K$.
\end{definition}

This is equivalent to compactness in the subspace topology. Indeed,
every set open in $K$ has the form $K\cap V$ with $V$ open in $X$, so we
can pass between the two kinds of cover without changing whether a
finite subcover exists.

\begin{theorem}\label{thm:compact-closed}
Every compact subset of a Hausdorff space is closed.
\end{theorem}

\section{The proof, with its dependencies visible}
\begin{proof}
Let $X$ be a Hausdorff space and let $K\subseteq X$ be compact. If
$K=\varnothing$, then $K$ is closed because its complement is the open
set $X$. Hence we may assume that $K$ is nonempty.

To prove that $K$ is closed, it suffices to prove that $X\setminus K$ is
open. Fix an arbitrary point $x\in X\setminus K$. We will construct an
open neighborhood $U$ of $x$ contained in $X\setminus K$.

For each $y\in K$, we have $x\ne y$. Since $X$ is Hausdorff, there are
open sets $U_y,V_y\subseteq X$ satisfying
\begin{equation}\label{eq:pair-separation}
  x\in U_y,\qquad y\in V_y,\qquad U_y\cap V_y=\varnothing.
\end{equation}
The family $\set{V_y}_{y\in K}$ is an open cover of $K$: each point
$y\in K$ belongs to its own set $V_y$. By compactness and
Definition~\ref{def:compact}, there are finitely many points
$y_1,\ldots,y_n\in K$ such that
\begin{equation}\label{eq:finite-cover}
  K\subseteq V_{y_1}\cup\cdots\cup V_{y_n}.
\end{equation}
Because $K$ is nonempty, the finite subcover may be taken with $n\geq1$.
Define
\begin{equation}\label{eq:chosen-neighborhood}
  U=\bigcap_{i=1}^{n}U_{y_i},
  \qquad
  V=\bigcup_{i=1}^{n}V_{y_i}.
\end{equation}
The set $U$ is open because it is a finite intersection of open sets.
Also $x\in U$, since $x\in U_{y_i}$ for every $i$. Thus $U$ is an
open neighborhood of $x$. Equation~\ref{eq:finite-cover} gives
$K\subseteq V$.

We claim that $U\cap V=\varnothing$. Otherwise, choose $z\in U\cap V$.
Since $z\in V$, there is an index $j\in\set{1,\ldots,n}$ with
$z\in V_{y_j}$. Since $z\in U$, we also have $z\in U_{y_j}$.
Consequently $z\in U_{y_j}\cap V_{y_j}$, contradicting
Equation~\ref{eq:pair-separation}. This proves the claim.

As $K\subseteq V$ and $U\cap V=\varnothing$, we obtain
$U\cap K=\varnothing$, or equivalently $U\subseteq X\setminus K$.
Every point $x\in X\setminus K$ therefore has an open neighborhood
contained in $X\setminus K$. The complement is the union of those open
neighborhoods, so it is open. If the complement is empty, it is open by
the topology axioms. In either case $K$ is closed.
\end{proof}

\section{Why the finite subcover matters}
It is tempting to set $U=\bigcap_{y\in K}U_y$ before using compactness.
This intersection would contain $x$ and avoid $K$, but it need not be
open. A topology guarantees openness only for finite intersections.
Equation~\ref{eq:finite-cover} is precisely what permits the finite
intersection in Equation~\ref{eq:chosen-neighborhood}.

The two hypotheses have distinct jobs. The Hausdorff property supplies
the pairwise separations in Equation~\ref{eq:pair-separation}.
Compactness supplies the finite selection. Finally, the topology axioms
assemble that finite selection into one usable neighborhood.

\section{A hypothesis we cannot discard}
\begin{example}[A compact set that is not closed]
Let $X=\set{a,b}$ have the indiscrete topology
$\set{\varnothing,X}$. The singleton $K=\set{a}$ is compact: every open
cover of $K$ contains $X$, and that one set is a finite subcover. But
$K$ is not closed, because $X\setminus K=\set{b}$ is not open. The
space is not Hausdorff, so this does not contradict
Theorem~\ref{thm:compact-closed}.
\end{example}

\begin{corollary}
If $K$ is compact in a Hausdorff space, then $\closure{K}=K$.
\end{corollary}
\begin{proof}
The closure $\closure{K}$ is the smallest closed subset containing $K$.
Theorem~\ref{thm:compact-closed} says that $K$ is already closed, so
$\closure{K}\subseteq K$. The reverse inclusion holds for the closure
of every subset. Therefore equality follows.
\end{proof}

\begin{exercise}
Read the construction of $U$ and $V$ again. Explain why it proves the
stronger statement that a point outside a compact subset of a Hausdorff
space can be separated from that subset by disjoint open sets.
\end{exercise}
\end{document}
